\chapter{Document Classes, Packages, and Page Design}
\label{chap:design}
\index{document class}\index{package}\index{page design}\index{geometry package}\index{header}

Once content is structured, design can support it consistently. The class
should make the largest decisions; packages should solve defined problems; and
local formatting should be the last, smallest layer. Reversing that order
produces a preamble full of repairs whose interactions are difficult to trace.

\begin{learningoutcomes}
  \item Choose among the standard \code{article}, \code{report}, and
        \code{book} classes.
  \item Apply class options and load packages for stated purposes.
  \item Locate and read current package documentation.
  \item Set page size, margins, spacing, headers, footers, and numbering.
  \item Prepare title, landscape, and multi-column pages when justified.
  \item Explain how \pkg{microtype} improves typography without manual repair.
\end{learningoutcomes}

\section{Choose the broad architecture}

The standard classes share many commands but differ in assumptions:

\begin{description}[style=nextline,leftmargin=1.6em]
  \item[article] Short scholarly works. Sections are the highest normal
  division; there is no \cmd{chapter}.
  \item[report] Longer works with chapters, a title page, and report-like
  front matter.
  \item[book] Books and long theses that need front, main, and back matter,
  two-sided conventions, and parts or chapters.
\end{description}

Choose from structure, not title. A document called ``Annual Report'' may be a
six-page article-class file; a doctoral thesis may use \code{book} or an
institutional class.

\begin{pausepredict}
What happens if source containing \cmd{chapter} is compiled with
\code{article}? Predict whether the class knows that command. Then decide
whether changing the class or changing the structure better matches the task.
\end{pausepredict}

\section{Class options}

Options in square brackets modify supported broad choices:

\begin{latexcode}{Class options}
\documentclass[11pt,a4paper,twoside,openright]{book}
\end{latexcode}

The example requests an 11-point base size, A4 paper, two-sided layout, and
chapters beginning on right-hand pages. Options are class-specific. An unknown
option may generate a warning or be passed to another component. Read the class
documentation rather than assuming an option exists.

For a publisher class, preserve its class line and options unless the current
instructions say otherwise. A generic draft and a publisher submission are
different design contexts.

\section{Packages solve named problems}

Load a package in the preamble with \cmd{usepackage}. Several packages can be
listed together only when they need no separate options and grouping remains
readable:

\begin{latexcode}{Purposeful package loading}
\usepackage{amsmath}  % Structured mathematics
\usepackage{graphicx} % External graphics
\usepackage{booktabs} % Professional table rules
\end{latexcode}

Comments record purpose. Package order sometimes matters: \pkg{hyperref} is
usually loaded late and \pkg{cleveref} after it. A class may already load a
package or provide an alternative interface.

\begin{commonerror}
\textbf{Symptom:} ``Option clash for package ...''.

\textbf{Cause:} the same package was loaded more than once with incompatible
options, perhaps once by the class.

\textbf{Repair:} search the preamble and log, load the package once with the
needed options, or pass class-compatible options before it is loaded. Test a
minimal example before changing a publisher template.
\end{commonerror}

\section{Read documentation at the point of need}

On a local TeX installation, the terminal command

\begin{terminalcode}{Open a package manual}
texdoc geometry
\end{terminalcode}

asks the TeX distribution for the \pkg{geometry} documentation. CTAN package
pages link to manuals and stable package information. Check the manual version
against the installed package when behaviour matters.

Read the introduction, examples, option reference, and compatibility notes.
Copying an isolated answer from the web may omit its class, package version, or
required loading order.

\section{Page size and margins}

The class establishes defaults. When a generic report has explicit page
requirements, \pkg{geometry} states them clearly:

\begin{latexcode}{Explicit generic report geometry}
\usepackage[
  a4paper,
  inner=30mm,
  outer=25mm,
  top=25mm,
  bottom=25mm,
  headheight=15pt
]{geometry}
\end{latexcode}

Inner and outer margins adapt to odd and even pages in a two-sided document.
Binding requirements can justify an asymmetric layout. Do not use geometry to
override a journal class that owns the submission dimensions.

Page design is a system: paper, margins, type size, line length, header, footer,
figures, and tables interact. Narrow margins may create long, tiring lines and
leave no safe binding area.

\section{Line spacing and paragraphs}

Some review or institutional workflows require double spacing. The
\pkg{setspace} package offers declared environments and commands:

\begin{latexcode}{Optional double-spaced draft}
\usepackage{setspace}

\begin{document}
\doublespacing
% Draft manuscript content
\end{document}
\end{latexcode}

Use this only when required. Increasing \cmd{baselinestretch} by guess can
affect footnotes and floats unexpectedly. Paragraph indentation and paragraph
spacing are also design choices; avoid combining a large indent with a large
blank gap unless a verified style calls for both.

\section{Headers, footers, and page numbering}

The \pkg{fancyhdr} package can define restrained running heads:

\begin{latexcode}{Two-sided running head pattern}
\usepackage{fancyhdr}
\pagestyle{fancy}
\fancyhf{}
\fancyhead[LE]{\thepage\quad Short Book Title}
\fancyhead[RO]{\nouppercase{\rightmark}\quad\thepage}
\renewcommand{\headrulewidth}{0.4pt}
\end{latexcode}

\cmd{fancyhf}\required{} clears existing fields. \code{LE} means left on even
pages; \code{RO} means right on odd pages. \cmd{rightmark} contains heading
information. Title and chapter-opening pages commonly use a separate
\code{plain} style.

\cmd{pagenumbering}\required{roman} and
\cmd{pagenumbering}\required{arabic} change numbering and reset the page
counter. Long documents should use structural commands such as \cmd{frontmatter}
and \cmd{mainmatter} when the class provides them.

\section{Title pages and special orientations}

The \code{titlepage} environment creates a page whose contents you control.
Use it for a report cover or institutionally defined thesis title page after
verifying required wording. Keep identity fields as commands or placeholders
so they are edited once.

The \pkg{pdflscape} package can rotate a genuinely wide page. The
\pkg{multicol} package can create balanced columns for a glossary or compact
reference, but multi-column prose complicates wide equations, figures, and
accessibility. A two-column journal class should control its own layout.

\begin{scienceinpractice}
An environmental assessment may use portrait pages for narrative, one
landscape page for a wide comparison matrix, and an appendix for detailed site
records. Orientation expresses a content need; it should not vary for visual
novelty.
\end{scienceinpractice}

\section{Microtypography}

The \pkg{microtype} package applies subtle character protrusion and font
expansion where supported. These adjustments can improve justification and
reduce distracting spacing without changing the argument or asking the author
to insert manual spaces.

\begin{latexcode}{A restrained typographic improvement}
\usepackage{microtype}
\end{latexcode}

It is intentionally uneventful. Good microtypography is often noticed as
comfortable reading rather than as an effect. Inspect the build log and output;
do not add discretionary hyphens throughout source merely to silence one draft
line.

\begin{goodpractice}
Start from the class defaults. Add one package when a named requirement appears.
Record the purpose, compile, inspect, and remove settings that no longer solve
a problem. A shorter preamble is easier to test and transfer.
\end{goodpractice}

\chapterproject

\begin{miniproject}
Create two build variants of the running article by changing only documented
settings: a readable generic article and a review draft with line spacing or
margins required by a hypothetical instruction. Keep the content, labels, and
bibliography identical. Record the variant purpose in the README and do not
present the hypothetical format as a universal submission rule.
\end{miniproject}

\chaptersummary

Classes provide broad architecture; options modify supported choices; packages
add named capabilities. Current manuals are part of the workflow. Geometry,
spacing, headers, footers, numbering, title pages, orientation, and columns form
an interacting page system. \pkg{microtype} can improve fine typography.
Publisher and institutional classes should retain control over prescribed
design.

\chaptercheck
\begin{chapterchecklist}
  \item The class matches the document's structural scale.
  \item Every package has a stated purpose and compatible loading order.
  \item Page settings follow verified requirements or documented preferences.
  \item Headers and page numbers work on odd, even, and opening pages.
  \item Special orientations and columns solve real content problems.
  \item The preamble contains no unexplained copied block.
\end{chapterchecklist}

\chapterexercisestitle
\begin{chapterexercises}
  \item Choose a standard class for a short paper, a laboratory report with
        chapters, and a textbook.
  \item Explain four options in the sample \code{book} declaration.
  \item Use \code{texdoc} to locate one installed package manual.
  \item Set asymmetric inner and outer margins for a two-sided generic report.
  \item Create a running header and a separate plain opening-page style.
  \item Diagnose a package option clash in a minimal example.
  \item Explain why a publisher class should not be restyled casually.
  \item Challenge: remove unnecessary packages from a copied preamble and
        document the evidence for each removal.
\end{chapterexercises}
